\section{微积分基本定理：变化与累积互相撤销}
\label{sec:02-Calculus/04-Integration/02}

上一节把定积分定义为 Riemann 和的极限。按那个定义去算 \(\int_1^2 x^2\,dx\)，流程如下：把 \([1,2]\) 等分成 \(n\) 段，\(\Delta x=1/n\)，右端点 \(x_i^*=1+i/n\)。Riemann 和展开：

\[
\sum_{i=1}^n \bigl(1+\frac{i}{n}\bigr)^2\cdot\frac{1}{n}
= \frac{1}{n}\sum_{i=1}^n \bigl(1+\frac{2i}{n}+\frac{i^2}{n^2}\bigr)
= \frac{1}{n}\Bigl(n+\frac{2}{n}\cdot\frac{n(n+1)}{2}+\frac{1}{n^2}\cdot\frac{n(n+1)(2n+1)}{6}\Bigr).
\]

整理、化简、取 \(n\to\infty\)——最终得到 \(7/3\)。算一个 \(x^2\) 就动用了平方求和公式。换成 \(\int_0^{\pi/2}\sin(x^2)\,dx\)？没有已知的闭合求和公式，直接走 Riemann 和这条路是死胡同。

但现在注意一个细节：\(x^3/3\) 的导数恰好是 \(x^2\)。如果我们只需要找到一个其导数是 \(x^2\) 的函数，代入上下限做差——\(2^3/3-1^3/3=8/3-1/3=7/3\)。这和 Riemann 和的苦力活算出的结果一模一样，但操作从''求一堆乘积之和的极限''退化成了''猜一个函数再代两个数''。

这种轻松不是特例。导数表和积分表之间有一股暗流——求导和积分像是在互相撤销。但这股暗流的精确形式是什么？它有哪些前提？它有没有失效的时候？

\subsection{累积到当前位置：定义一个会生长的函数}

固定起点 \(a\)，把上限变成一个可以滑动的变量 \(x\)：

\[
G(x)=\int_a^x f(t)\,dt.
\]

\(G(x)\) 的读法：从起点 \(a\) 走到当前位置 \(x\)，\(f\) 总共累积了多少。给一个新的 \(x\)，就吐出一个新的累积总量——\(G\) 是个函数，不是一个数。

现在问一个关键问题：\(G(x)\) 对 \(x\) 的变化率是多少？换个问法——把上限从 \(x\) 再往前推一小步 \(h\)，累积量会多出来多少？

\[
G(x+h)-G(x)=\int_x^{x+h} f(t)\,dt.
\]

这个新增量就是 \(f\) 在小区间 \([x,x+h]\) 上的积分。除以区间长度 \(h\)：

\[
\frac{G(x+h)-G(x)}{h} = \frac{1}{h}\int_x^{x+h} f(t)\,dt.
\]

等号右边是 \(f\) 在这段小区间上的平均值——积分值除以宽度。到此还没用到任何深奥的东西，只是把差商重新包装成了平均值的形式。

\subsection{连续让平均值回到原函数}

如果 \(f\) 在 \(x\) 处连续，那么 \(h\) 足够小时，区间 \([x,x+h]\) 上的每一个 \(f(t)\) 都离 \(f(x)\) 很近——近到它们在平均值里造成的偏差可以小到任意预设的程度。严格写出来：

任给 \(\varepsilon>0\)，由连续性存在 \(\delta>0\)，使得当 \(|t-x|<\delta\) 时 \(|f(t)-f(x)|<\varepsilon\)。取 \(|h|<\delta\)（以 \(h>0\) 为例），区间上每一点都满足这个不等式，于是

\[
\begin{aligned}
\Bigl|\frac{1}{h}\int_x^{x+h} f(t)\,dt - f(x)\Bigr|
&= \Bigl|\frac{1}{h}\int_x^{x+h} \bigl[f(t)-f(x)\bigr]\,dt\Bigr| \\[4pt]
&\le \frac{1}{h}\int_x^{x+h} \bigl|f(t)-f(x)\bigr|\,dt \\[4pt]
&< \frac{1}{h}\int_x^{x+h} \varepsilon\,dt = \varepsilon.
\end{aligned}
\]

差商与 \(f(x)\) 的距离可以小于任意正数——这意味着差商的极限就是 \(f(x)\)。即

\[
G'(x)=f(x).
\]

这是微积分基本定理的第一部分，常被概括为一句话：先积再导，回到原函数。累积操作被求导操作撤销了。连续函数的累积函数，其瞬时变化率恰好等于被积函数在当前点的取值。

\subsection{从第一部分导出第二部分：原函数让定积分变成减法}

第一部分告诉我们 \(G(x)=\int_a^x f(t)\,dt\) 满足 \(G'=f\)。假设我们另有某个函数 \(F\)，也满足 \(F'=f\)。那么 \(G\) 和 \(F\) 的导数在整个区间上相等，它们的差 \(G-F\) 的导数恒为零——中值定理保证导数恒为零的函数必为常数。因此存在 \(C\)，使

\[
G(x)-F(x)=C,\quad \text{对所有 }x.
\]

代入 \(x=a\)：\(G(a)=\int_a^a f=0\)，所以 \(0-F(a)=C\)，即 \(C=-F(a)\)。于是

\[
G(x)=F(x)-F(a).
\]

取 \(x=b\)：

\[
\int_a^b f(t)\,dt = G(b) = F(b)-F(a).
\]

这就是第二部分——Newton-Leibniz 公式。计算定积分，不再需要切碎、采样、求和、取极限。找一个其导数是 \(f\) 的函数 \(F\)，把上下限代入相减即可。

检验 \(x^2\) 在 \([1,2]\)：\(F(x)=x^3/3\)，\(F(2)=8/3\)，\(F(1)=1/3\)，差 \(7/3\)。注意方括号记号 \(\bigl[F(x)\bigr]_a^b=F(b)-F(a)\)——它是一步纯减法的缩写，不是把上下限直接塞进被积函数。

\subsection{净变化定理：积分无处不在的通用面孔}

第二部分换一种读法，适用范围急剧扩大。若 \(Q'(t)=r(t)\)，则

\[
Q(b)-Q(a)=\int_a^b r(t)\,dt.
\]

一句话：存量的变化等于变化率的累积。骨架干净到只剩符号，但它的具体化身铺满应用领域：

\begin{itemize}
\tightlist
\item 水箱水量变化 = 净流入速率对时间的积分。进水 3 升/分钟，出水 1 升/分钟，净速率 2 升/分钟，10 分钟内水量增加 \(\int_0^{10} 2\,dt=20\) 升。
\item 位置变化（位移）= 速度对时间的积分。第一秒速度从 0 匀加速到 10 米/秒，位移 = \(\int_0^1 10t\,dt=5\) 米。
\item 速度变化 = 加速度对时间的积分。恒定加速度 \(9.8\) 米/秒\(^2\) 持续 3 秒，速度增加 \(\int_0^3 9.8\,dt=29.4\) 米/秒。
\item 概率在区间上的累积 = 概率密度在该区间的积分。
\end{itemize}

这些不是类比——它们是同一个数学结构在不同语境下被不同的物理名目称呼而已。如果你还知道初值 \(Q(a)=Q_0\)，任何时刻的存量都可以复原：

\[
Q(t)=Q_0+\int_a^t r(s)\,ds.
\]

这是后续微分方程积分形式的原型——给定变化规则和初始状态，就能重建全部轨迹。例如已知物体从静止开始运动，加速度 \(a(t)=2t\)（匀变加速度），初速度 \(v(0)=0\)。速度变化由加速度积分给出：\(v(t)=0+\int_0^t 2s\,ds=t^2\)。位置由速度积分给出：\(x(t)=x_0+\int_0^t s^2\,ds=x_0+t^3/3\)。两层积分嵌套，每次都在取消一次求导——微积分基本定理的叠加版本。

\subsection{为什么这个定理是''基本''的}

微积分基本定理之所以冠名''基本''，不是因为它简单——从 Riemann 和定义到差商均值到极限论证，这条链并不短——而是因为它把两个看似不相干的操作焊接在了一起。微分是局部的、点上的、着眼于变化的；积分是全局的、区间上的、着眼于累积的。基本定理说：它们是互逆的。

这一互逆关系是后续所有积分技巧的基石。下一节的换元积分和分部积分，就是从这条''互逆''出发，把链式法则和乘积法则反向翻译成积分语言。如果没有基本定理架起的桥，这些翻译就无处落脚。

\subsection{不定积分和定积分：共用一个符号的不同对象}

初学者最容易被三样东西的相似外表搞混：

\begin{itemize}
\tightlist
\item 不定积分 \(\int f(x)\,dx = F(x)+C\)：原函数族。那个 \(+C\) 不能丢——常数的导数为零，中值定理保证了在一个区间上任意两个原函数至多相差常数。
\item 定积分 \(\int_a^b f(x)\,dx\)：一个数。计算时常数在端点相减抵消，所以从不写 \(+C\)。
\item 变上限积分 \(\int_a^x f(t)\,dt\)：一个函数，对上限求导回到被积函数。
\end{itemize}

三种东西共用 \(\int\) 符号，但对象的数学类型截然不同。在计算定积分时常数之所以自然消去，不是''运算法则规定可以省略''——是端点代入做差时它代数地被消掉了。

积分变量是内绑哑元，叫什么名字不改变积分值：

\[
\int_a^x f(t)\,dt = \int_a^x f(u)\,du.
\]

用 \(t\) 而非 \(x\) 做内部变量，只为避免积分上限 \(x\) 与被积表达式里的变量名撞车——是记号卫生，不带深层差异。

\subsection{连续性条件不是写在纸上的空话}

基本定理的标准陈述以''设 \(f\) 连续''开头。这不是套话。如果被积函数在某点跳跃，累积函数在该点连续但不可导——出现一个折角。

考虑符号函数在 \([-1,1]\) 上：\(x<0\) 时值为 \(-1\)，\(x>0\) 时值为 \(1\)。从 \(-1\) 到 \(x\) 的累积是

\[
\int_{-1}^x \sign(t)\,dt = |x|-1,\quad x\in[-1,1].
\]

累积函数 \(|x|-1\) 在 \(x=0\) 处有一个尖角——连续但不可导。被积函数在 \(x=0\) 的值甚至没有明确定义，累积函数的导数也因而在 \(x=0\) 不存在。''先积再导回到原函数''在跳变点失效——定理的''如果''精确地划出了安全的领地，领地之外需要更强的积分理论。

更一般的条件下（Lebesgue 积分、绝对连续函数）基本定理有更强的版本，但入门阶段抓住这条分界线就够了：\(f\) 连续，累积函数光滑且导数完美回到 \(f\)；\(f\) 跳跃，累积函数出折角，不能处处求导回到原函数。定理的精确性在于告诉你安全范围，不是模糊地说''一般都对''。
